% Oversells «Sheʼriyat» shabloni (XeLaTeX yoki LuaLaTeX).
% Har bir sheʼr: \section*{Nomi}, keyin \begin{uslub}{sher} … \end{uslub};
% qatorlar \\ bilan, bandlar orasida boʻsh qator. Koʻrinish sahifa.sty da.
\documentclass[11pt]{article}
\usepackage[a5paper]{geometry}
\usepackage{fontspec}
\setmainfont{Noto Serif}
\usepackage{babel}
\babelprovide[import=uz-Latn, main]{uzbek}
\def\uzbekhyphenmins{22}
\usepackage{sahifa} % Oversells sahifa dizayni (sahifa.sty)

\begin{document}

\begin{center}
  \vspace*{20mm}
  {\Huge\bfseries Bahor nafasi\par}
  \vspace{8mm}
  {\large Sheʼrlar\par}
\end{center}
\clearpage

\section*{Bahor}
\begin{uslub}{sher}
Bogʻlarda oppoq oʻrik gullari,\\
Shamolda tebranar tolning shoxlari.\\
Qaldirgʻoch qaytibdi uzoq yoʻllardan,\\
Uyimiz ustida uning uyasi.

Ariqda shildirab oqadi suvlar,\\
Dalada koʻkarar bugʻdoy maysalar.\\
Har tong yangi umid olib keladi —\\
Bahor nafasini sezar har kishi.
\end{uslub}
\begin{uslub}{sana}
Toshkent, 2026-yil, mart
\end{uslub}

\section*{Ona tilim}
\begin{uslub}{sher}
Onam aytgan allada\\
Ilk bor eshitdim seni,\\
\matnuslubi{urgu}{Ona tilim}, sen bilan\\
Anglab yetdim dunyoni.
\end{uslub}

\end{document}
